We use the lazy symmetric normalization of \citet[Section~2, equation~(3)]{fountoulakis2022open}.

\subsection{Graph and optimization problems}
\label{subsec:shared-source-aligned-problem}

We use boldface for vectors and matrices and regard vectors as columns.
Let \(\mathcal{G}=(\mathcal{V},\mathcal{E})\) be a finite connected simple
unweighted undirected graph with \(\mathcal{V}=[n]:=\{1,\ldots,n\}\),
\(n\geq2\), and adjacency matrix \(\bm{A}\in\{0,1\}^{n\times n}\).
For \(i\in\mathcal{V}\), define
\[
    \mathcal N(i):=\{j\in\mathcal{V}:\{i,j\}\in\mathcal{E}\},
    \qquad
    d_i:=|\mathcal N(i)|,
    \qquad
    \bm{D}:=\diag(d_1,\ldots,d_n).
\]
Degrees always refer to the original graph, including in restricted
systems. Connectedness gives \(d_i\geq1\), so \(\bm{D}^{-1/2}\) is well
defined. For \(\mathcal{S}\subseteq\mathcal{V}\), its neighborhood,
external boundary, and volume are
\begin{equation*}
    \mathcal N(\mathcal{S}):=\bigcup_{i\in\mathcal{S}}\mathcal N(i),
    \qquad
    \partial\mathcal{S}:=\mathcal N(\mathcal{S})\setminus\mathcal{S},
    \qquad
    \vol(\mathcal{S}):=\sum_{i\in\mathcal{S}}d_i.
\end{equation*}
Write \(\supp(\bm{x}):=\{i\in\mathcal{V}:x_i\neq0\}\).
We use \(\preceq,\succeq\) for Loewner order and \(\leq,\geq\) for
entrywise inequalities between vectors or matrices. The seed distribution
satisfies
\begin{equation*}
    \bm{s}\in\R_+^n,
    \qquad
    \inner{\one}{\bm{s}}=1.
\end{equation*}
Here \(\one\) is the all-ones vector. Our main results use a point source
\(\bm{s}=\bm{e}_v\), where \(\bm{e}_v\) is the \(v\)-th standard basis
vector; general seeds and their input and work costs are treated in
\cref{app:seed-scope}. For a teleportation parameter \(\alpha\in(0,1]\),
define
\begin{equation}
\begin{aligned}
    \bm{\mathcal L}&:=\bm{I}-\bm{D}^{-1/2}\bm{A}\bm{D}^{-1/2},\\
    \bm{Q}&:=\alpha\bm{I}+\frac{1-\alpha}{2}\bm{\mathcal L}
      =\frac{1+\alpha}{2}\bm{I}
       -\frac{1-\alpha}{2}\bm{D}^{-1/2}\bm{A}\bm{D}^{-1/2},\\
    \bm{b}&:=\alpha\bm{D}^{-1/2}\bm{s}.
\end{aligned}
    \label{eq:shared-pagerank-matrices}
\end{equation}
The normalized Laplacian satisfies
\(\bm{0}\preceq\bm{\mathcal L}\preceq2\bm{I}\), hence
\(\alpha\bm{I}\preceq\bm{Q}\preceq\bm{I}\). The PageRank quadratic is
\begin{equation}
    f(\bm{x}):=\frac12\inner{\bm{x}}{\bm{Q}\bm{x}}-\inner{\bm{b}}{\bm{x}},
    \qquad
    \nabla f(\bm{x})=\bm{Q}\bm{x}-\bm{b}.
    \label{eq:shared-pagerank-objective}
\end{equation}
Its unique minimizer and the corresponding lazy personalized PageRank
(PPR) vector are
\begin{equation*}
    \bm{x}_0^*:=\bm{Q}^{-1}\bm{b},
    \qquad
    \bm{\pi}:=\bm{D}^{1/2}\bm{x}_0^*.
\end{equation*}
For \(\rho>0\), the regularized PageRank (RPPR) objective and its unique
minimizer are defined by
\begin{equation}
    g_\rho(\bm{x}):=\alpha\rho\norm{\bm{D}^{1/2}\bm{x}}_1,
    \qquad
    F_\rho(\bm{x}):=f(\bm{x})+g_\rho(\bm{x}),
    \qquad
    \bm{x}_\rho^*:=\argmin_{\bm{x}\in\R^n}F_\rho(\bm{x}).
    \label{eq:shared-rppr-objective}
\end{equation}
We abbreviate the minimizer as \(\bm{x}^*\) when the objective is clear.
Writing \(\mathcal{S}_\rho^*:=\supp(\bm{x}_\rho^*)\), standard RPPR
properties give \(\bm{x}_\rho^*\geq\bm{0}\),
\(\vol(\mathcal{S}_\rho^*)\leq1/\rho\), and the coordinatewise
optimality conditions
\begin{equation}
    \nabla_i f(\bm{x}_\rho^*)
    \in
    \begin{cases}
        \{-\alpha\rho\sqrt{d_i}\}, & (\bm{x}_\rho^*)_i>0,\\
        [-\alpha\rho\sqrt{d_i},0], & (\bm{x}_\rho^*)_i=0.
    \end{cases}
    \label{eq:shared-rppr-kkt}
\end{equation}
For a point source $\vs=\eunit{v}$, $\vx_\rho^*\neq\vzero$ exactly when
\begin{equation}
 0<\rho<1/d_v.
 \label{eq:nonzero-rppr-regime}
\end{equation}
Indeed, the optimality condition at zero is
$\vb\leq\alpha\rho\mD^{1/2}\one$, which reduces to $\rho\geq1/d_v$.
A single degree query identifies the zero case. When $\alpha=1$, the
solution is explicit:
$\vx_\rho^*=((1-\rho d_v)_+/\sqrt{d_v})\eunit{v}$, where
$(t)_+:=\max\{t,0\}$.

\subsection{Accuracy guarantees}

The RPPR task takes $\rho>0$ and an additive tolerance
$\epsobj>0$, and asks for a sparse vector $\widehat\vx$ satisfying
\begin{equation}
 F_\rho(\widehat\vx)-F_\rho(\vx_\rho^*)\leq\epsobj.
 \label{eq:rppr-output-target}
\end{equation}
The output lists the nonzero coordinates of $\widehat\vx$ and their values;
unlisted coordinates are zero. For PPR, an output $\widehat\vx$ represents
$\widehat\vpi=\mD^{1/2}\widehat\vx$. The semantic target is
\begin{equation}
 \norm{\mD^{-1}(\widehat\vpi-\vpi)}_\infty
 =\norm{\mD^{-1/2}(\widehat\vx-\vx_0^*)}_\infty
 \leq\epsppr.
 \label{eq:semantic-ppr-target}
\end{equation}
Thus $\epsppr$ measures solution error per unit degree. The PageRank matrix $\mQ$ is a \emph{Stieltjes matrix}: symmetric positive
definite with nonpositive off-diagonal entries.
With $\vomega=\mD^{1/2}\one$, the identities
\begin{equation*}
 \mQ\vomega=\alpha\vomega,\qquad
 \mQ^{-1}\geq\vzero,
 \qquad |\mQ|\vomega=\vomega
\end{equation*}
will be used throughout; $|\mQ|$ denotes entrywise absolute values.
Since $\mA\one=\mD\one$, we have
$\mD^{-1/2}\mA\mD^{-1/2}\vomega=\vomega$, which gives the first identity
and, using the zero diagonal of $\mA$, the third. Also,
$\mI-\mQ\geq\vzero$ entrywise and
$\norm{\mI-\mQ}_2\leq1-\alpha<1$, so the Neumann series
$\mQ^{-1}=\sum_{k=0}^{\infty}(\mI-\mQ)^k$ converges and is entrywise
nonnegative.

\begin{proposition}[Regularization bias and objective-to-PPR conversion]
\label{prop:rppr-ppr-bridge}
For every $\rho>0$,
\begin{equation}
 \vzero\leq\vx_0^*-\vx_\rho^*\leq\rho\vomega.
 \label{eq:rppr-bias}
\end{equation}
Every nonnegative output satisfying \cref{eq:rppr-output-target} obeys
\begin{equation}
 \norm{\mD^{-1/2}(\widehat\vx-\vx_0^*)}_\infty
 \leq\rho+\sqrt{2\epsobj/\alpha}.
 \label{eq:objective-to-ppr}
\end{equation}
\end{proposition}
\begin{proof}
The KKT conditions give
$\vzero\leq\vb-\mQ\vx_\rho^*\leq\alpha\rho\vomega$.
Multiply by $\mQ^{-1}\geq\vzero$ and use
$\mQ^{-1}\vomega=\vomega/\alpha$ to obtain \cref{eq:rppr-bias}.
Strong convexity gives
$\norm{\widehat\vx-\vx_\rho^*}_2\leq\sqrt{2\epsobj/\alpha}$.
Since $d_i\geq1$, its degree-normalized maximum norm is no larger.
The triangle inequality proves \cref{eq:objective-to-ppr}.
\end{proof}
For example, to achieve the semantic PPR target \cref{eq:semantic-ppr-target}, it suffices to set
\begin{equation}
 \rho=\epsppr/2,\qquad \epsobj=\alpha\epsppr^2/8
 \label{eq:direct-ppr-budgets}
\end{equation}
is sufficient for \cref{eq:semantic-ppr-target}. No identification of the
two tolerances is involved. A second, residual-based sufficient condition is
\begin{equation}
 \norm{\mD^{-1/2}(\vb-\mQ\widehat\vx)}_\infty\leq\alpha\epsppr
 \quad\Longrightarrow\quad
 \norm{\mD^{-1/2}(\widehat\vx-\vx_0^*)}_\infty\leq\epsppr.
 \label{eq:ppr-residual-certificate}
\end{equation}
It follows from inverse positivity applied to the two coordinatewise
residual bounds. Our main RPPR algorithm uses its proved objective budget;
it does not need to scan the whole graph to check this sufficient PPR
certificate.

\subsection{Local access and computational cost}

The algorithms receive the seed label $v$, their numerical parameters, and
degree and adjacency-list access to $\mathcal G$. All other vertex labels
must be discovered through incident entries. No support information or
auxiliary graph data are supplied, and both algorithms use no graph-wide
preprocessing. Unless stated otherwise, we use an exact-real algebraic word model and
define the \emph{total work} of a complete execution by
\begin{equation}
 \Work:=N_{\mathrm{adj}}+N_{\mathrm{deg}}+N_{\mathrm{op}}+N_{\mathrm{out}}.
 \label{eq:total-work}
\end{equation}
Here $N_{\mathrm{adj}}$ counts adjacency-list entry inspections,
$N_{\mathrm{deg}}$ degree queries, $N_{\mathrm{op}}$ all remaining scalar
arithmetic operations, comparisons, random choices, and state reads or
writes, and $N_{\mathrm{out}}$ the words written to return the output. Each such operation
costs one unit. Composite operations, including dictionary updates,
local matrix and solver construction, threshold reporting, certification,
and cleanup, are charged for all their elementary operations.

Every inspection is counted, including repeated access to cached
incidences. Scanning a set $\mathcal S$ reads only its vertices' lists and
contributes $\vol(\mathcal S)$ to $N_{\mathrm{adj}}$. Full scans therefore
contribute $\sum_j d_{u_j}$, where $u_j$ is the vertex scanned at the
$j$th scan and repeated visits are included. For randomized algorithms,
$\Work$ includes failed trials and restarts; expected bounds control
$\E[\Work]$. Peak storage is reported separately.

An exact scalar occupies one word regardless of its encoding length.
For rational inputs, a density $f_i$ and original degree $d_i$ represent
$x_i=f_i\sqrt{d_i}$ and $\pi_i=d_if_i$, avoiding square-root computations.
The deterministic schedule uses dyadic halving, arithmetic, and comparisons.
\Cref{app:bounded-arithmetic} separately charges numerical and label
encoding costs for a specified bounded-arithmetic implementation;
floating-point stability requires separate analysis.

In soft bounds, $\widetilde{\mathcal{O}}$ hides only polylogarithms in inverse numerical
parameters, requested accuracy, explicitly supplied failure probability, and
the number of exposed records. It hides no ambient graph-size factor, no
unknown minimum sign margin, and no uncharged setup. In the nonzero regime
\cref{eq:nonzero-rppr-regime}, set
\begin{equation}
 L_{\mathrm{par}}:=\log\frac{2}{\alpha\rho\min\{1,\epsobj\}}.
 \label{eq:main-parameter-log}
\end{equation}

\subsection{Main results}

The next two theorems establish the RPPR guarantee, and
\cref{cor:main-ppr} gives its PPR consequence.
A nonnegative vector $\widehat\vx$ satisfying $\mQ\widehat\vx\leq\vb$
is called a \emph{subsolution}. Both algorithms can return subsolutions
lying coordinatewise below $\vx_\rho^*$ at the stated work bounds.

\begin{theorem}[Deterministic local acceleration]
\label{thm:deterministic-rppr}
For every graph and point source above, $0<\alpha\leq1$,
$0<\rho<1/d_v$, and $\epsobj>0$, \cref{alg:deterministic-rppr} returns
\begin{equation*}
 \vzero\leq\widehat\vx\leq\vx_\rho^*,\qquad
 F_\rho(\widehat\vx)-F_\rho(\vx_\rho^*)\leq\epsobj,
 \qquad\vol(\supp\widehat\vx)\leq1/\rho.
\end{equation*}
Its total work, defined in \cref{eq:total-work}, satisfies
\begin{equation}
 \Work=\mathcal{O}\!\left(\frac{L_{\mathrm{par}}}{\rho\sqrt\alpha}
           \log^2\!\left(2+\frac{L_{\mathrm{par}}}{\rho\sqrt\alpha}\right)\right)
 =\softO\!\left(\frac1{\rho\sqrt\alpha}\right),
 \label{eq:deterministic-main-work}
\end{equation}
and its peak storage is $\mathcal{O}(1+L_{\mathrm{par}}/(\rho\sqrt\alpha))$ words.
For $\rho\geq1/d_v$ it returns the exact zero solution in constant word work.
\end{theorem}

The proof occupies \cref{sec:deterministic-method}. The bound improves
the graph-independent worst-case work from
$\softO(1/(\alpha\rho))$ to $\softO(1/(\rho\sqrt\alpha))$.
The principal new estimate controls the cumulative kinetic volume of an
entire stage while allowing accelerated updates outside the optimal support.

\begin{theorem}[Randomized local acceleration, with support adaptation]
\label{thm:randomized-rppr-main}
Let $k_*=|\mathcal S_\rho^*|$ and $V_*=\vol(\mathcal S_\rho^*)$.
For $0<\alpha\leq1$, $0<\rho<1/d_v$, and $\epsobj>0$, there is a Las
Vegas local algorithm returning $\vzero\leq\widehat\vx\leq\vx_\rho^*$
satisfying \cref{eq:rppr-output-target}. It can also guarantee
$\vzero\leq\vb-\mQ\widehat\vx\leq2\alpha\rho\vomega$.
Its expected total work satisfies
\begin{equation}
 \E[\Work]=\softO\!\left(V_*\min\{k_*,\alpha^{-1/2}\}\right)
 \leq\softO\!\left(
       \min\left\{\frac1{\rho^2},\frac1{\rho\sqrt\alpha}\right\}\right).
 \label{eq:randomized-adaptive-work}
\end{equation}
Alternatively, a single capped run reports an explicit failure with
probability at most $\zeta\in(0,1)$, has the same expected work up to
failure-probability logarithms, and always satisfies the output guarantee
when it returns a vector. All support admissions are deterministically
certified.
\end{theorem}

The factors suppressed in \cref{eq:randomized-adaptive-work} include
logarithmic objective-accuracy dependence in the batch count and in the SDD
solves. The refined bound follows because every nonempty batch adds a new
vertex, while \cref{thm:threshold-batch-depth} bounds the number of batches
spectrally. Thus the randomized method retains the small-support bound as
well as the accelerated worst-case envelope: up to logarithmic factors,
the smaller of support size and the inverse square root of the teleportation
parameter determines how many restricted systems need to be solved.
Its graph and boundary records have size $\mathcal{O}(V_*)$; numerical workspace is
that of one SDD call on the current active set. Allocation is charged, so expected peak
space is in particular bounded by the displayed expected work.

\begin{corollary}[Local PPR at semantic accuracy]
\label{cor:main-ppr}
For $0<\epsppr<1$, degree-normalized PPR error $\epsppr$ can be attained
in $\softO(1/(\epsppr\sqrt\alpha))$ deterministic work, or in the same
expected work by a Las Vegas algorithm. The deterministic output obeys
$\vzero\leq\widehat\vx\leq\vx_0^*$, and both outputs have support volume
at most $2/\epsppr$. If $\epsppr\geq1$, zero suffices.
\end{corollary}
\begin{proof}
Use \cref{eq:direct-ppr-budgets} and
\cref{prop:rppr-ppr-bridge}. The logarithm of the requested objective
accuracy remains a sum of $\log(1/\alpha)$ and $\log(1/\epsppr)$.
For $\epsppr\geq1$, use $\vpi\geq\vzero$, $\one^\trans\vpi=1$, and
$d_i\geq1$.
\end{proof}

\paragraph{Bounded arithmetic.}
For rational numerical inputs, \cref{thm:bounded-deterministic-rppr} gives
an implementation with the same $\softO(1/(\rho\sqrt\alpha))$ local
operation bound using integers of controlled length. With schoolbook
integer arithmetic its bit cost is
$\softO(B^2/(\rho\sqrt\alpha))$, where $B$ explicitly includes the
input numerical encodings, encountered label and degree encodings, and
parameter and precision logarithms. This statement has no hidden dependence
on a minimum activation margin. It concerns a specified rounded algorithm,
not an arbitrary numerical implementation of the exact recurrence.
